\section{Related Work}
\label{sec:related_work}

\para{Chain-of-thought prompting.}
\citet{wei2022chain} showed that prompting LLMs to produce intermediate reasoning steps dramatically improves performance on arithmetic, commonsense, and symbolic reasoning tasks.
Subsequent work has refined CoT through self-consistency decoding~\citep{wang2023selfconsistency}, automatic prompt construction, and task decomposition.
These methods place all reasoning \emph{before} the answer, treating deliberation as a preamble rather than an ongoing process.
\iscotshort differs by interleaving reasoning \emph{between} sentences of the visible output, so that each sentence is preceded by its own deliberation step.

\para{Pause and thinking tokens.}
\citet{goyal2024think} introduced learnable \texttt{<pause>} tokens that, when appended to the input during both pretraining and finetuning, improve downstream task accuracy on benchmarks such as SQuAD (+18\% EM) and CommonSenseQA (+8\%).
\citet{herel2024thinking} proposed inserting thinking tokens after each word in LSTM-based language models, observing per-sentence perplexity improvements on reasoning-heavy sentences but slight degradation on average perplexity.
\citet{galashov2025catch} made pause tokens adaptive, allowing the model to request additional computation by emitting \texttt{<DON'T KNOW>} tokens when uncertain.
All three approaches require training modifications.
We test whether the same intuition---giving the model more time to think---can be realized purely through prompting, using the model's native instruction-following ability.

\para{Latent reasoning.}
\citet{hao2024coconut} introduced Coconut, which feeds the model's last hidden state back as the next input embedding to enable reasoning in a continuous latent space.
\citet{lovelace2025star} integrated latent diffusion planning into autoregressive generation, pausing generation to perform diffusion-based planning in sentence embedding space before resuming.
These methods achieve sentence-level planning but require architectural modifications and specialized training.
Our work tests whether a similar effect---sentence-level deliberation---can be achieved with an off-the-shelf model through structured prompting alone.

\para{Test-time compute scaling.}
\citet{snell2024scaling} systematically studied test-time compute allocation, showing that compute-optimal strategies per prompt can improve efficiency by 4$\times$ over best-of-$N$ sampling.
\citet{muennighoff2025s1} demonstrated that appending ``Wait'' tokens to extend a model's thinking process leads to self-correction and improved performance on competition mathematics, with their s1-32B model exceeding o1-preview on MATH and AIME24 benchmarks.
These studies focus on maximizing benchmark accuracy through more thinking.
We complement this line of work by studying how additional test-time compute changes the \emph{qualitative character} of the text, not just its correctness.

\para{LLM-as-judge evaluation.}
Our evaluation relies partly on LLM-as-judge pairwise comparisons~\citep{zheng2023judging}.
This approach is known to exhibit systematic biases, particularly toward longer and more verbose responses.
We observe this bias directly in our results and discuss its implications for evaluating methods that trade breadth for per-sentence quality.
